
\documentclass{article}
\usepackage{colortbl}
\usepackage{makecell}
\usepackage{multirow}
\usepackage{supertabular}

\begin{document}

\newcounter{utterance}

\centering \large Interaction Transcript for game `ifeval', experiment `ifeval', episode 288 with selfrefine{-}llama{-}3.1{-}8b{-}instruct.
\vspace{24pt}

{ \footnotesize  \setcounter{utterance}{1}
\setlength{\tabcolsep}{0pt}
\begin{supertabular}{c@{$\;$}|p{.15\linewidth}@{}p{.15\linewidth}p{.15\linewidth}p{.15\linewidth}p{.15\linewidth}p{.15\linewidth}}
    \# & \multicolumn{2}{c}{Player} && \multicolumn{2}{c}{Game Master} \\
    \hline

    \theutterance \stepcounter{utterance}  
    & & & \multicolumn{4}{p{0.6\linewidth}}{
        \cellcolor[rgb]{0.9,0.9,0.9}{
            \makecell[{{p{\linewidth}}}]{
                \texttt{\tiny{[P1$\langle$GM]}}
                \texttt{Give me a logical question about chatting and conquering. Make it interesting. Highlight at least one section of your answer with markdown, i.e. *highlighted section*. Put your entire response inside double quotation marks.} \\
            }
        }
    }
     \\ \\

    \theutterance \stepcounter{utterance}  
    & \multicolumn{4}{p{0.6\linewidth}}{
        \cellcolor[rgb]{0.561,0.737,0.561}{
            \makecell[{{p{\linewidth}}}]{
                \texttt{\tiny{[P1$\rangle$GM]}}
                \texttt{"Imagine a medieval kingdom where the ruler has the power to conquer neighboring lands through a unique form of diplomacy: a high{-}stakes chat game. The ruler can challenge other kingdoms to a chat battle, where the goal is to outwit and outmaneuver the opponent through clever wordplay and strategic thinking.} \\
\\ 
\texttt{The twist is that the ruler has a special ability: the power to *steal* a set number of words from the opponent's vocabulary, rendering them unable to use those words in the chat battle. The ruler can use this power strategically to disrupt the opponent's plans and gain an advantage.} \\
\\ 
\texttt{The question is: if the ruler has a vocabulary of 1000 words and can steal up to 20 words from the opponent's vocabulary per turn, what is the maximum number of turns it would take for the ruler to steal all 1000 words from the opponent, assuming the opponent has a vocabulary of 2000 words and uses a fixed number of words in each turn, but the ruler does not know what that number is?"} \\
            }
        }
    }
    & & \\ \\

    \theutterance \stepcounter{utterance}  
    & & & \multicolumn{2}{p{0.3\linewidth}}{
        \cellcolor[rgb]{0.9,0.9,0.9}{
            \makecell[{{p{\linewidth}}}]{
                \texttt{\tiny{[GM$|$GM]}}
                {'detectable_format:number_highlighted_sections': {'num_highlights': 1}, 'startend:quotation': {}}
            }
        }
    }
    & & \\ \\

    \theutterance \stepcounter{utterance}  
    & & & \multicolumn{2}{p{0.3\linewidth}}{
        \cellcolor[rgb]{0.9,0.9,0.9}{
            \makecell[{{p{\linewidth}}}]{
                \texttt{\tiny{[GM$|$GM]}}
                \texttt{game\_result = WIN} \\
            }
        }
    }
    & & \\ \\

\end{supertabular}
}

\end{document}
